% TOP_PROBLEM: {"id": 30006939, "title": "Balanced L4 boundedness of the trilinear Hilbert transform", "category_id": 8, "category": {"id": 8, "name": "analysis", "display_name": "Analysis", "description": "Limits, continuity, calculus, and function theory.", "slug": "analysis", "order_index": 8, "created_at": "2026-07-31T15:26:25.671Z"}, "status": "open"}
% FIRST_POSTED: 2026-09-25
% ENTRY_KIND: statement_only
% !TeX program = lualatex
% Definition and sources only; this file is not an LLM solution attempt.
\documentclass[11pt]{article}
\usepackage[margin=25mm]{geometry}
\usepackage{fontspec}
\setmainfont{Latin Modern Roman}
\usepackage{amsmath,amssymb,mathtools}
\usepackage{unicode-math}
\setmathfont{Latin Modern Math}
\usepackage{xurl,hyperref}
\hypersetup{colorlinks=true,urlcolor=blue}
\setcounter{secnumdepth}{0}
\setlength{\parindent}{0pt}
\setlength{\parskip}{5pt}
\setlength{\emergencystretch}{3em}
\begin{document}
\section*{\texorpdfstring{Balanced L4 boundedness of the trilinear Hilbert transform}{Balanced L4 boundedness of the trilinear Hilbert transform}}\label{30006939}
\subsection{Definitions and mathematical statement}
For Schwartz functions $f,g,h:{\mathbb{R}}\to{\mathbb{C}}$, define the trilinear Hilbert transform with fixed slopes $1,2,3$ by
\[
 T(f,g,h)(x)=\operatorname{p.v.}\int_{{\mathbb{R}}}
 f(x+t)g(x+2t)h(x+3t)\,\frac{{\,\mathrm{d}} t}{t}.
\]
Principal value means symmetric deletion of an interval around zero. The selected boundedness question is the existence of an absolute finite $C$ with
\[
 \|T(f,g,h)\|_{L^{4/3}({\mathbb{R}})}
 \le C\|f\|_{L^4({\mathbb{R}})}\|g\|_{L^4({\mathbb{R}})}\|h\|_{L^4({\mathbb{R}})}.
\]
The estimate is required uniformly over all such inputs and would give a bounded extension to the indicated spaces. The exponent identity $3/4=1/4+1/4+1/4$ fixes the balanced scaling; this entry does not ask for the entire possible exponent range.
\par\smallskip\textit{Formulation and concept references:} \href{https://terrytao.wordpress.com/2007/05/10/open-question-boundedness-of-the-trilinear-hilbert-transform/}{[S1]}.

\subsection{Short English statement}
Is the trilinear Hilbert transform bounded from three copies of L-four into L-four-thirds at the specified slopes?
\subsection{Sources}
\begin{itemize}
\item[S1]Terence Tao, Open question: boundedness of the trilinear Hilbert transform (10 May 2007).
\url{https://terrytao.wordpress.com/2007/05/10/open-question-boundedness-of-the-trilinear-hilbert-transform/}.
\textit{Formulation source.}
\end{itemize}
\end{document}
